\scanpage{184}
\mspage{II.32}
% archive sheet number: 183
% block: top, lines 1-2 | text
formées avec les $G_\alpha$, qu'on peut considérer comme modules à action analytique.

% block: upper third, lines 3-4 | text
6.4. Le syst. projectif des $G_\alpha$ définit un système projectif d'Algèbres de Lie $\mathfrak g_\alpha$ sur $\mathbb Q_\ell$

% block: upper third, line 5 | display
\[
\deleted{\unsure{\mathfrak g(\pi)=\mathfrak g_\ell(\pi)=(\mathfrak g_\alpha)}}\qquad
\mathfrak g=\operatorname{Lie}^{\mathrm{an}}(\pi)=(\mathfrak g_\alpha)
\]

% block: upper third, line 6 | text
et on a un foncteur

% block: upper to middle third, lines 7-11 | diagram
\diagram{\begin{tikzcd}[ampersand replacement=\&,column sep=large]
\& \deleted{\operatorname{Mod}(\unsure{G_\alpha})} \& \\
\operatorname{Mod}_{\mathbb Q_\ell}(\pi) \arrow[ur,no head,"{\unsure{\simeq}}"] \arrow[rr] \arrow[d]
\& \& \operatorname{Mod}\mathfrak g_{\deleted{\unsure{\ell}}} \arrow[d,no head] \\
\begin{gathered}\text{catégorie des}\\\text{\unsure{les modules|modules} $\ell$-adiques}\\\text{sur $\pi$ (sur $\mathbb Q_\ell$)}\end{gathered}
\& \& \begin{gathered}\varinjlim\text{ des}\\\text{catégories de}\\\text{modules }\deleted{\text{sur les}}\mathfrak g_\alpha\\\text{de dim finie}\\\text{sur }\mathbb Q_\ell\end{gathered}
\end{tikzcd}}

% block: middle third, lines 12-16 | text
Ce foncteur est fidèle, mais en général \uline{pas} pleinement fidèle: On a perdu en passant de \inserted{$\unsure{\pi^{\mathrm{an}}}$} $G$ à $\mathfrak g$ (\unsure{intuitivement}, on perd ``des composantes connexes). Mais soit $(U_i)$ le syst. fond. des \unsure{s-groupes|sous groupes} ouverts de $\pi$. Alors l'inclusion $U_i\mathbin{-}\pi$ définit un \uline{isom} de pro-Lie

% block: lower third, line 17 | display
\[
\mathfrak g(U_i)\simeq\mathfrak g(\pi)
\]

% block: lower third, line 18 | text
d'où un foncteur

% block: lower third, line 19 | display
\[
\operatorname{Mod}_{\mathbb Q_\ell}(U_i)\longrightarrow\operatorname{Mod}\mathfrak g
\]

% block: lower third, line 20 | text
et par passage à la limite

% block: lower third, line 21 | display
\[
\boxed{\varinjlim_i\operatorname{Mod}_{\mathbb Q_\ell}(U_i)\longrightarrow\operatorname{Mod}\mathfrak g}
\]

% block: bottom, line 22 | text
Cette flèche est une équivalence de \uline{catégories}
